\chapter{Introduction}

\section{Background and Motivation}

Modern grocery and convenience retail is transitioning rapidly from passive digital search and spreadsheet-based replenishment toward \textbf{agentic AI}---software agents capable of perceiving multimodal inputs (voice, text, images, and transactional time series), invoking structured tools, reasoning under business constraints, and collaborating safely with human decision-makers. Unlike static FAQ chatbots or offline forecasting notebooks, an \textbf{AI agent} operates inside an explicit observe--reason--act loop: it parses ambiguous natural-language intent, queries deterministic ground-truth databases or analytical engines, enforces domain safety and policy constraints, and returns verifiable actions or recommendations.

Globally, leading grocery chains have demonstrated both the commercial upside and the engineering difficulty of deploying retail AI agents across 2024--2026:
\begin{itemize}
    \item \textbf{Walmart (\emph{Sparky}):} Deployed an agentic shopping assistant capable of converting high-level shopper goals (such as meal or event planning under a budget) and photos of handwritten lists into concrete carts, item substitutions, and cross-platform commerce workflows, reporting rapid consumer adoption and higher basket sizes~\cite{walmart_sparky}.
    \item \textbf{Amazon (\emph{Rufus}):} Embedded a generative shopping agent trained on catalog attributes, reviews, and Q\&A data to answer complex comparative queries across voice, text, and visual inputs, extending toward autonomous purchasing flows while relying on proprietary ecosystem scale~\cite{amazon_rufus}.
    \item \textbf{Carrefour (\emph{Hopla} / \emph{Hopla+}):} Pioneered European grocery conversational shopping by linking dietary objectives (lower sugar, Nutri-Score optimization, recipe-to-basket assembly) directly to supermarket inventory, while simultaneously deploying internal enterprise copilots for store operations~\cite{carrefour_hopla}.
\end{itemize}

In \textbf{Vietnam}, modern retail is undergoing one of the fastest structural modernizations in Southeast Asia, yet operational friction remains acute across thousands of neighborhood supermarkets and convenience stores:
\begin{itemize}
    \item \textbf{Thanh Lam Mart (Thanh Lam Mart / Thanh Lam Mart):} Operating over 4,500--5,000 retail points nationwide, Thanh Lam Mart treats data and machine learning as the backbone of its consumer operating system to automate centralized ordering, reduce manual inventory counting, and free up to 30\% of store-staff time for customer service~\cite{vneconomy_retail_ai}.
    \item \textbf{Bách Hóa Xanh (MWG):} Competing heavily in fresh food and fast-moving perishables, Bách Hóa Xanh relies on centralized data-driven distribution and short-horizon inventory control to minimize perishable spoilage (\emph{shrinkage}) while keeping shelves freshly stocked~\cite{bhx_supplychain}.
    \item \textbf{FPT Retail (Long Châu, FPT Shop) \& FPT.AI:} Combines internal AI-assisted supply-chain optimization with Vietnamese voicebots (\emph{Voicebot CSAT}) and omnichannel ERP/CRM-connected chatbots to handle high-volume customer lookup and routine inquiries 24/7~\cite{fpt_ai_retail}.
\end{itemize}

Despite this momentum, a critical gap separates enterprise headquarters analytics from \textbf{store-level execution in Vietnam}. On the \textbf{sales floor}, shoppers at physical supermarkets rarely have access to a fast, noise-resilient Vietnamese voice assistant that accurately compares unit prices, promotions, sugar levels, and multi-item meal budgets across thousands of local and international brands (e.g., \emph{Thanh Lam Home}, \emph{TH True Milk}, \emph{ACEFOODS}, \emph{OMO}) without hallucinating prices, stock levels, or medical claims. Meanwhile, in the \textbf{back office}, store managers at independent or branch-level convenience stores still juggle daily purchase orders, lot-level expiration dates (1--3 day perishables), sudden demand events, and supplier returns manually, often without transparent demand recovery for out-of-stock days or safe human-in-the-loop (HITL) guardrails.

This capstone project directly addresses both sides of the retail store within one unified dual-agent architecture implemented across two repositories: \texttt{Retail-AI-} (the customer-facing Vietnamese voice kiosk agent, \emph{Mai}) and \texttt{ops\_chest} (\url{https://github.com/quanchu14104/ops_chest.git}, the manager-facing perishable store operations agent).

\section{Literature Review}

Building trustworthy retail AI agents requires synthesizing four distinct research and engineering threads: (i) tool-using LLM agents and safety guardrails, (ii) real-time speech enhancement and streaming ASR in noisy environments, (iii) perishable inventory optimization and profit-aligned machine learning forecasting, and (iv) industry retail agent deployments.

\textbf{LLM AI Agents, Grounded RAG, and Guardrails.}
Yao et al.~\cite{yao2023react} introduced \emph{ReAct}, demonstrating that interleaving reasoning traces with external tool calls reduces hallucination compared to pure chain-of-thought generation. Schick et al.~\cite{schick2023toolformer} and Lewis et al.~\cite{lewis2020rag} established that delegating arithmetic, retrieval, and factual lookup to external tools and indexed knowledge stores dramatically improves factual precision. However, unconstrained agent loops introduce two severe production failure modes: \emph{latency amplification} (each sequential planner/selector/reviewer LLM round adds 2--4\,s of dead wait) and \emph{safety miscalibration}. Röttger et al.~\cite{rottger2024xstest} showed via \emph{XSTest} that naive safety filters suffer high false-refusal rates on benign domain queries, while the OWASP Top 10 for LLM Applications~\cite{owasp_llm} cataloged prompt injection, sensitive data disclosure, and excessive agency as primary operational threats. Consequently, production retail agents require deterministic pre-LLM gates, strict tool-grounded output verification, and zero direct-execution rights for high-impact operational mutations.

\textbf{Real-Time Speech Enhancement and Streaming ASR.}
Physical supermarket kiosks operate amid 5--15\,dB ambient noise from HVAC hum (50/100\,Hz), trolley clatter, and multi-talker babble. Classical spectral subtraction degrades tonal Vietnamese syllables, whereas modern neural speech enhancement models such as \emph{DeepFilterNet}~\cite{schroter2022deepfilternet} and its lightweight Dual-Path Deep Filtering descendant \texttt{DPDFNet2}~\cite{sherpa_onnx} suppress non-stationary background noise on CPU with low distortion. Paired with recurrent Voice Activity Detection (\emph{Silero VAD}~\cite{silero_vad}) and streaming cloud ASR supporting domain phrase biasing (\emph{Azure AI Speech} \texttt{PhraseListGrammar}~\cite{azure_speech}), voice pipelines can process audio concurrently while the user speaks instead of waiting for post-utterance file uploads. Nevertheless, real-world kiosk deployments must also solve loudspeaker self-echo during full-duplex barge-in and Vietnamese-accented English brand slips (e.g., \emph{Thanh Lam Home} misheard as \emph{Thanh Lam Mart hôn}).

\textbf{Perishable Retail Forecasting and Inventory Optimization.}
Retail inventory literature highlights two classical pitfalls in fresh-food replenishment. First, raw point-of-sale records are \emph{right-censored}: whenever a shelf stocks out mid-day, recorded sales underestimate true demand, causing naive forecasters to enter a downward stockout spiral~\cite{nahmias2011perishable}. Recent open benchmarks such as \emph{FreshRetailNet-50K}~\cite{freshretailnet50k} expose hourly stockout flags across real convenience stores, enabling empirical demand recovery and quantile forecasting with gradient-boosted trees (\emph{LightGBM}~\cite{ke2017lightgbm}). Second, optimizing statistical accuracy (such as WAPE) does not automatically maximize store profit on short-shelf-life goods, because asymmetric overage (spoilage) and underage (lost margin and lost traffic) costs are governed by the perishable \emph{Newsvendor} critical fractile~\cite{silver2017inventory} and First-Expired-First-Out (FEFO) lot depletion.

\textbf{Comparative Synthesis and Research Gap.}
Table~\ref{tab:lit_comparison} contrasts representative academic methods and commercial retail AI deployments against the requirements of Vietnamese retail chains.

\begin{table}[H]
\centering
\caption{Comparison of reviewed research approaches and retail AI industry systems}
\label{tab:lit_comparison}
\footnotesize
\begin{tabularx}{\textwidth}{L{2.9cm} L{2.7cm} L{4.2cm} Y}
\toprule
\textbf{Work / System} & \textbf{Approach} & \textbf{Advantages} & \textbf{Limitations for Our Problem} \\
\midrule
ReAct~\cite{yao2023react} \& RAG~\cite{lewis2020rag} & Multi-step LLM reasoning + external tool retrieval & Flexible handling of open-ended multi-turn queries; reduces parametric hallucination & High multi-call latency (5--14\,s); lacks retail unit-price, safety, and HITL execution boundaries \\
Walmart Sparky~\cite{walmart_sparky} \& Amazon Rufus~\cite{amazon_rufus} & Proprietary omnichannel shopping agents & High conversion, visual/list-to-cart flows, deep review/order integration & Closed cloud ecosystems built for US/English mobile apps; no Vietnamese noisy-kiosk voice or store-floor ops \\
Carrefour Hopla~\cite{carrefour_hopla} & Grocery dietary chatbot + recipe-to-basket & Strong nutritional/budget grocery curation & Text/web-centric; separate from lot-level FEFO expiry and store-manager replenishment \\
Thanh Lam Mart / BHX / FPT.AI~\cite{vneconomy_retail_ai,fpt_ai_retail} & Centralized ML ordering + call-center voicebots & Proven scale across Vietnamese chains; cuts manual stocktaking time & Headquarters analytics and call-center bots remain decoupled from in-store voice advisory and transparent manager HITL chat \\
DeepFilterNet~\cite{schroter2022deepfilternet} \& Silero VAD~\cite{silero_vad} & Local neural denoising + recurrent speech detection & Real-time CPU noise suppression and accurate speech boundary detection & Standalone audio modules; require integration with streaming ASR, brand repair, and anti-echo barge-in \\
LightGBM~\cite{ke2017lightgbm} on FreshRetailNet~\cite{freshretailnet50k} & Gradient-boosted quantile demand forecasting & Captures nonlinear lags, weekday/holiday seasonality, and stockout features & Minimizing WAPE alone can over-order and hurt store profit on small or intermittent SKUs without profit-gated routing \\
\bottomrule
\end{tabularx}
\end{table}

\textbf{Targeted Research and Engineering Gap.}
As summarized in Table~\ref{tab:lit_comparison}, existing systems solve fragments of the retail problem in isolation. The concrete gap targeted by this capstone is designing a \textbf{verifiable, low-latency dual-agent retail architecture tailored to Vietnam} that: (1) eliminates unnecessary multi-call planner bottlenecks on the sales floor while preserving zero-hallucination catalog grounding and clinical/OWASP safety; (2) sustains accurate Vietnamese-plus-English-brand voice recognition under 12\,dB--5\,dB supermarket noise with sub-second post-speech latency; and (3) pairs profit-gated LightGBM/EWMA perishable replenishment and FEFO expiry analytics with strict runtime numeric grounding and HITL manager governance.

\section{Problem Statement}

At the store level of a Vietnamese supermarket or convenience chain, two interdependent problem surfaces degrade service quality and operating margin simultaneously:

\begin{enumerate}
    \item \textbf{Sales-Floor Customer Advisory Problem (\texttt{Retail-AI-}):} Shoppers regularly ask natural, noisy Vietnamese questions spanning brand filtering (\emph{``khăn mặt cotton Thanh Lam Home''}), dietary constraints (\emph{``sữa tươi không đường''}), unit-price value (\emph{``nước giặt OMO loại nào tiết kiệm nhất?''}), two-brand comparisons, and budget-capped meal baskets (\emph{``mua đồ nấu lẩu thái dưới 300k''}). Naive LLM architectures fail on four concrete axes: (a) \emph{multi-call planner latency} (2--4 sequential non-streaming LLM calls taking 5.4--14.2\,s before the first token appears); (b) \emph{noisy speech degradation and upload delay} (post-utterance WAV upload + cloud Whisper takes 2.0--3.0\,s and exhibits 14.2--22.2\% WER under supermarket noise, while naive catalog post-correction corrupts everyday Vietnamese words); (c) \emph{half-duplex speaker echo} preventing natural interruption; and (d) \emph{hallucination and safety risk} on unverified prices, stock availability, allergens, or medical claims.
    \item \textbf{Back-Office Store Operations Problem (\texttt{ops\_chest}):} Store managers overseeing hundreds of perishable SKUs with 1--3 day shelf lives must decide daily order quantities, markdown discounts, supplier returns, and write-offs. Heuristic ordering suffers 14.5\% overall stockouts, 8.7\% stockouts on traffic-driving items, and 5.3\% waste, while raw sales underestimate true demand by $-32.3\%$ on stockout days. Conversely, blindly applying ML forecasts across every SKU segment over-orders on small stores (degrading median profit by $-5.6\%$), and giving an LLM chatbot unconstrained write permissions creates severe operational risk if the model invents figures or claims to have dispatched purchase orders without approval.
\end{enumerate}

The core research question is therefore: \emph{How can we engineer a unified dual-agent AI architecture that solves sales-floor Vietnamese voice consultation in real time with zero factual hallucination, while equipping store managers with profit-validated perishable inventory intelligence under strict Human-in-the-Loop governance?}

\section{Objectives and Contributions}

\subsection*{General Objective}
Design, implement, and empirically benchmark a dual-agent AI platform for Vietnamese retail store chains uniting a real-time, noise-resilient voice kiosk assistant (\texttt{Retail-AI-}) and a profit-gated, HITL-governed perishable store operations agent (\texttt{ops\_chest}).

\subsection*{Specific Objectives}
\begin{itemize}
    \item \textbf{Build a low-latency, noise-resilient Vietnamese voice pipeline (\texttt{Retail-AI-}):} Integrate local Silero VAD~v6, real-time WebSocket streaming STT (\code{/audio/transcribe/stream}) backed by Azure AI Speech SDK with Thanh Lam Mart \texttt{PhraseListGrammar}, local \texttt{DPDFNet2} ONNX neural speech enhancement with SNR-adaptive clean bypass and heavy-noise rescue, catalog-anchored English-brand repair, full-duplex three-layer anti-echo barge-in, and progressive clause-level Azure Neural TTS.
    \item \textbf{Eliminate unnecessary planner latency via grounded SSE streaming (\texttt{Retail-AI-}):} Architect a modular zero-planner streaming engine (\code{streaming_consultation.py}) plus one-call slim basket planner that emits verified catalog cards within tens of milliseconds and streams live DeepSeek tokens across common shopping queries, while routing clinical, allergy, infant, and OWASP red-team queries through deterministic pre-LLM safety gates and bounded verification.
    \item \textbf{Deliver profit-gated perishable forecasting and HITL store operations (\texttt{ops\_chest}):} Implement hourly censored-demand recovery, 29-feature LightGBM quantile forecasting (P10/P50/P90), a cross-validated profit-gated ML/EWMA forecast router, perishable Newsvendor and FEFO lot-level replenishment simulation, expiry disposition workflows, and a 31-tool manager chatbot guarded by runtime numeric verification and mandatory human approval.
    \item \textbf{Define a fail-closed cross-module data contract and run rigorous empirical benchmarks:} Specify the read-only branch-snapshot contract uniting \texttt{ops\_chest} and \texttt{Retail-AI-}, and validate both repositories via automated test suites, controlled 100-question SOTA~V2 evaluations, multi-SNR speech ablations, 10-query live streaming audits, and 50-store out-of-sample profit backtests.
\end{itemize}

\subsection*{Contributions}
\begin{itemize}
    \item \textbf{Theoretical and methodological contributions:} Demonstrates three transferable agentic-AI engineering principles: (i) \emph{tiered zero-planner Groq LPU streaming with Turn-1 semantic caching and Singapore (\texttt{sin1}) edge co-location}---separating fast conversational streaming from complex structured JSON safety planning cuts live HTTPS latency from 2.4--5.5\,s down to \textbf{508--870\,ms} cold and \textbf{125--140\,ms} cached (0 tokens) without sacrificing factual grounding; (ii) \emph{multi-source three-tier catalog knowledge enrichment with an 8-stage zero-defect cleaning pipeline} (inspired by Hugging Face \code{datatrove} and OpenFoodFacts) expanding verified back-of-pack ingredient coverage to 1,214+18 SKUs while failing closed on unverified SKUs; and (iii) \emph{profit-gated forecast routing}---selecting ML versus statistical forecasts by simulated store profit and coupling LightGBM point medians with EWMA dispersion outperforms pure WAPE-based ML selection out-of-sample.
    \item \textbf{Practical contributions:} Delivers two runnable, heavily tested open-architecture prototypes tailored to Vietnamese retail chains: a live-deployed mobile-first Vietnamese voice and visual kiosk assistant on Vercel Singapore (\url{https://mai-retail-kiosk.vercel.app}) featuring 0.0\% brand WER under 12\,dB supermarket noise, a 99.9\%-coverage clickable Product Photo Lightbox, and real-time Zalo Bot manager escalation; and a containerized Docker store-operations copilot calibrated on real perishable demand that improves simulated store gross profit by $+4.9\%$ while preserving full human accountability.
\end{itemize}

\section{Scope and Limitations}

\subsection*{Scope}
This project implements and evaluates both modules at a pre-commercial engineering level: \texttt{Retail-AI-} runs both locally (\code{127.0.0.1:8000} / \code{5173}) and in live Vercel Production in Singapore (\texttt{sin1}, \url{https://mai-retail-kiosk.vercel.app}) over the 3,099-SKU Thanh Lam Mart catalog snapshot (enriched with 2,137 zero-defect product knowledge records), mobile QR UI with clickable product photo popups, live Zalo Bot manager alerting, and sandbox cart/checkout flow; \texttt{ops\_chest} runs as a FastAPI + React/Vite + PostgreSQL + Mailpit Docker stack calibrated on FreshRetailNet-50K real demand (Demo Store~18 and 50 held-out backtest stores) plus a 263-SKU synthetic business-rule test world.

\subsection*{Limitations}
\begin{itemize}
    \item \textbf{Static catalog snapshot and simulated retail transactions:} \texttt{Retail-AI-} operates on a frozen Thanh Lam Mart snapshot (15 September 2026) with clearly labeled simulated 30-day demo sales (\texttt{SIMULATED}) and sandbox checkout; it is not connected to live Thanh Lam Mart POS or ERP stock feeds.
    \item \textbf{Unbridged runtime synchronization between repositories:} Although both modules belong to one unified logical architecture with an explicit freshness-TTL branch contract (Appendix~C), live runtime synchronization from \texttt{ops\_chest}'s lot ledger into \texttt{Retail-AI-}'s kiosk snapshot is not yet deployed across a live store network.
    \item \textbf{Remaining structured-JSON cloud latency, partial nutrition coverage, and synthetic-to-field acoustic gap:} Although tiered Groq LPU streaming + semantic caching + \texttt{sin1} eliminated the earlier 1.2--1.5\,s streaming TTFT floor, complex structured JSON safety queries still rely on cloud DeepSeek endpoints, and $\approx 59\%$ of catalog SKUs still lack public back-of-pack nutrition/calorie labels (handled honestly as \texttt{UNKNOWN}); noise-resilience evaluations used controlled 12\,dB and 5\,dB supermarket babble/HVAC mixtures rather than a multi-week physical supermarket floor trial.
    \item \textbf{Assumed operational parameters in \texttt{ops\_chest}:} While daily demand in FreshRetailNet-50K is real, product display names, unit margins (62--75\% cost ratio), 1--3 day shelf lives, supplier schedules, and sandbox Mailpit PO dispatch are structured operational assumptions documented transparently in Section~3.1.
\end{itemize}
